\documentclass[../../../include/open-logic-section]{subfiles}

\begin{document}

\olfileid{sth}{ordinals}{replacement}
\olsection{Panggantian}

Ing \olref[sth][ordinals][vn]{sec}, kita nerangake
alesan ngenalake ordinal kanthi ngusulake supaya
ordinal dianggep minangka tipe urutan, yaiku wakil
kanonik tumrap urutan becik. Supaya cara iki lumaku,
kita kudu mbuktekake yen \emph{saben urutan becik
isomorfis karo sawenehing ordinal}. Iki ngidini
kita netepake $\ordtype{A, <}$ minangka ordinal
$\alpha$ kang nyukupi $\tuple{A, <} \isomorphic \alpha$.

Sayange, kita \emph{ora bisa} mbuktekake asil kang
dikarepake mung nganggo aksioma kang wis dikenalake
nganti saiki. (Sebabe bakal katon ing
\olref[replacement][strength]{sec}; saiki cukup
dimangerteni yen pancen ora bisa.) Kita butuh
gagasan anyar, yaiku:

\begin{axiom}[Skema Panggantian]
Kanggo sembarang rumus $\phi(x, y)$, iki sawijining aksioma:
\begin{quote}
	kanggo sembarang $A$, yen
	$(\forall x \in A)\lexists![y][\phi(x,y)]$, mula
	$\Setabs{y}{(\exists x \in A)\phi(x,y)}$ ana.
\end{quote}
\end{axiom}
\noindent
Kaya Pamisahan, iki sawijining skema: ngasilake
aksioma kang cacahé tanpa wates, siji kanggo saben
$\phi$ kang beda-beda. Skema iki uga bisa (lan
lumrahe pancen) ditulis mangkene:

\begin{defish}
Kanggo sembarang rumus $\phi(x,y)$ kang ora ngemot
``$B$'', iki sawijining aksioma:
\[
\forall A[(\forall x \in A)\lexists![y][\phi(x,y)] \lif \exists B\forall y (y \in B \liff (\exists x \in A)\phi(x,y))]
\]
\end{defish}

Nanging nalika sepisanan dideleng, rumus iki katon
ruwet. Akibat langsung saka Panggantian ing ngisor
iki mbokmenawa nuduhake gagasan intuitif kita kanthi
luwih \emph{cetha}: \footnote{Term iku ungkapan kang
nemtokake persis siji obyek (kanggo saben pangisian
variabel bebasa). Contone, ``$\emptyset$'' iku term
kang nemtokake himpunan kosong; ``$\{x\}$'' iku
term kang nemtokake himpunan tunggal saka $x$
(apa wae $x$); ``$x \cup y$'' iku term kang nemtokake
gabungan saka $x$ lan $y$ (apa wae nilai loro-lorone).}

\begin{cor}
Kanggo sembarang term $\tau(x)$ lan sembarang
himpunan $A$, himpunan iki ana:
\[
	\Setabs{\tau(x)}{x \in A} = \Setabs{y}{(\exists x \in A)y = \tau(x)}.
\]
\end{cor}

\begin{proof}
Amarga $\tau$ iku \emph{term},
$\forall x \lexists![y][\tau(x) = y]$. Mula luwih-luwih
$(\forall x \in A)\lexists![y][\tau(x) = y]$.
Dadi $\Setabs{y}{(\exists x \in A)\tau(x) = y}$
ana miturut Panggantian.
\end{proof}
\noindent
Mula ``Panggantian'' dadi jeneng kang trep kanggo
aksioma iki: saka sawijining himpunan $A$, kita
bisa mbentuk himpunan anyar,
$\Setabs{\tau(x)}{x \in A}$, kanthi ngganti saben
anggota $A$ nganggo bayangane miturut~$\tau$.
Miturut notasi kanggo bayangan himpunan dening
fungsi, kita bisa nulis $\funimage{\tau}{A}$
kanggo $\Setabs{\tau(x)}{x \in A}$.

Nanging kang wigati, $\tau$ iku \emph{term}.
Dheweke ora kudu dadi (jeneng kanggo) sawijining
\emph{fungsi} miturut \olref[sfr][fun][rel]{sec},
yaiku himpunan tartamtu saka pasangan urut.
Yen $f$ pancen fungsi miturut teges iku lan
ditetepake kanggo saben anggota himpunan kang
dirembug, himpunan
$\funimage{f}{A} = \Setabs{f(x)}{x \in A}$ mung
himpunan bagean tartamtu saka $\ran{f}$, lan
anane wis dijamin nganggo aksioma~$\Zminus$.
\footnote{Coba pikirna
$\Setabs{y \in \bigcup \bigcup f}{(\exists x \in A)y = f(x)}$.}
Kosok baline, Panggantian nambah kakuwatan
\emph{gedhe} marang aksioma kita, kaya kang bakal
katon ing \olref[sth][replacement][]{chap}.


\end{document}
